\hypersetup{pageanchor=false}
\pagestyle{empty}
\begin{titlepage}
\centering
\includegraphics[width=2.8cm]{figuras/logo_oficial.png}\par\medskip
{\large\bfseries Instituto Federal Goiano\\Campus Urutaí\par}
Curso de Sistemas de Informação\par
\vspace{1.8cm}
{\large Artur Duarte Monteiro\par}
\vspace{2.4cm}
{\LARGE\bfseries \titulo\par}
\vfill
Urutaí\\2026
\end{titlepage}
\hypersetup{pageanchor=true}
\setcounter{page}{1}
\thispagestyle{empty}
\begin{center}
{\large Artur Duarte Monteiro\par}
\vspace{2.5cm}
{\Large\bfseries \titulo\par}
\end{center}
\vspace{2cm}
\hfill\begin{minipage}{.53\textwidth}\singlespacing
Monografia desenvolvida no Curso de Sistemas de Informação do Instituto Federal Goiano, Campus Urutaí, como Trabalho de Curso.\par\medskip
Orientador: Gabriel da Silva Vieira.
\end{minipage}
\vfill
\begin{center}Urutaí\\2026\end{center}
\clearpage
\thispagestyle{empty}
\chapter*{Resumo}
\begin{singlespace}
A recuperação de fotografias de uma pessoa em coleções com múltiplos rostos exige distinguir identidade facial de semelhança visual geral. Este trabalho investiga se a incorporação de contexto visual global por fusão tardia melhora a recuperação baseada em identidade facial. Foi desenvolvido um protótipo local com detecção SCRFD, representação facial ArcFace por InsightFace e descritores globais de ResNet50 previamente treinada, sem novo treinamento. A fotografia inteira é a unidade recuperada; os escores faciais são agregados por máximo e combinados com a similaridade global mediante pesos fixos. A comparação principal usa 589 fotografias Gallagher, 20 consultas de 20 identidades, 19 fotografias-fonte e 588 candidatas por ranking. A relevância deriva das anotações, independentemente do sucesso do detector, e a fonte é excluída. Na execução de fechamento, a busca facial obteve precisão média das consultas (mAP) de 0,952261; a busca global obteve 0,280864; as fusões face/global 0,9/0,1, 0,7/0,3 e 0,5/0,5 obtiveram, respectivamente, 0,946274, 0,901460 e 0,813297. Nenhuma fusão avaliada superou a busca facial no mAP agregado. A fusão 0,9/0,1 apresentou dois ganhos, nove empates e nove perdas por consulta, mostrando que o resultado médio não implica ausência de ganhos locais. LFW e INRIA Holidays são verificações auxiliares com protocolos locais, e o álbum familiar participa apenas da aceitação funcional privada. A contribuição é uma comparação experimental delimitada e rastreável, apoiada em software funcional. Dependência entre consultas, seleção por frequência, pesos não otimizados e uma única coleção principal limitam a generalização. Não se demonstra causalidade do fundo nem significância estatística.
\par\medskip\noindent\textbf{Palavras-chave:} recuperação de imagens; identidade facial; contexto visual; fusão tardia; avaliação experimental.
\end{singlespace}
\clearpage
\chapter*{Abstract}
\begin{otherlanguage}{english}\begin{singlespace}
Retrieving photographs of a person from collections containing multiple faces requires distinguishing facial identity from overall visual similarity. This study investigates whether incorporating global visual context through late fusion improves identity-based photograph retrieval. A local prototype combines SCRFD detection, ArcFace representations through InsightFace, and global descriptors from a pretrained ResNet50, without additional training. Whole photographs are the retrieval unit; facial scores are aggregated by their maximum and combined with global similarity using fixed weights. The main comparison uses 589 Gallagher photographs, 20 queries representing 20 identities, 19 source photographs, and 588 candidates per ranking. Relevance is derived from annotations independently of detector success, and source photographs are excluded. In the final recorded execution, face-only retrieval achieved a mean average precision (mAP) of 0.952261, while global-only retrieval achieved 0.280864. Face/global fusion weights of 0.9/0.1, 0.7/0.3, and 0.5/0.5 achieved 0.946274, 0.901460, and 0.813297, respectively. None of the evaluated fusion settings exceeded the face-only aggregate mAP. The 0.9/0.1 setting improved two queries, tied nine, and degraded nine, demonstrating that the aggregate result does not imply the absence of individual gains. LFW and INRIA Holidays provide auxiliary checks under local protocols, while a family album is used only for private functional acceptance. The contribution is a bounded, traceable experimental comparison supported by functional software. Query dependence, frequency-based selection, unoptimized weights, and a single primary collection limit generalization. Neither background causality nor statistical significance is established.
\par\medskip\noindent\textbf{Keywords:} image retrieval; facial identity; visual context; late fusion; experimental evaluation.
\end{singlespace}\end{otherlanguage}
\clearpage
\listoffigures
\clearpage
\listoftables
\clearpage
\chapter*{Lista de abreviaturas e siglas}
\begin{description}
\item[AP] Average Precision (precisão média de uma consulta).
\item[CBIR] Content-Based Image Retrieval (recuperação de imagens por conteúdo).
\item[CNN] Convolutional Neural Network (rede neural convolucional).
\item[CPU] Central Processing Unit (unidade central de processamento).
\item[CUDA] Plataforma de computação paralela utilizada com GPU NVIDIA.
\item[GPU] Graphics Processing Unit (unidade de processamento gráfico).
\item[IA] Inteligência artificial.
\item[LFW] Labeled Faces in the Wild.
\item[LGPD] Lei Geral de Proteção de Dados Pessoais.
\item[mAP] Mean Average Precision (média das APs das consultas).
\item[ONNX] Open Neural Network Exchange.
\item[P@K] Precisão nas primeiras K posições.
\item[PIPA] People In Photo Albums.
\item[R@K] Revocação nas primeiras K posições.
\item[SCRFD] Sample and Computation Redistribution for Face Detection.
\item[SHA-256] Função de resumo criptográfico utilizada para verificar arquivos.
\item[TC] Trabalho de Curso.
\end{description}
\clearpage
\tableofcontents
\clearpage
